% sec-11-references-backmatter.tex - References and Back Matter (full)
% Unnumbered back matter so the numbered IND sections end at 11 (Relevant
% Information), matching the ReGARDD IND Table-of-Contents numbering.
\section*{References and Back Matter}
\addcontentsline{toc}{section}{References and Back Matter}

This back matter closes Initial Investigational New Drug Application v1.0 for the
LLM-Directed PDAC Robotic Daraxonrasib Phase 1 study, assembled within repository
v4.3.0. The complete reference list renders below through the \texttt{ieeetr}
bibliography style from \texttt{references.bib}. Every entry carries a clickable
URL and, where a digital object identifier is assigned, a clickable DOI URL with
the DOI text shown in full; long links break on any character so that no link runs
off the right margin of the text block. The reference base supporting this
application is intentionally comprehensive and spans seven evidentiary strands:
the FDA Phase 1 and combination drug-device regulatory guidances; the daraxonrasib
(RMC-6236) clinical and mechanistic literature underpinning the RASolute 302 and
RASolve 301 programs and the June 2025 Breakthrough Therapy designation; the
21 CFR Part 312 (IND) and Part 50 (informed consent) and ICH Physical Artificial
Intelligence regulatory adaptations; the five author computational works dated 2025
to 2026 (the 60-second PDAC Whipple plus daraxonrasib simulation, the artificial
intelligence digital-twin PDAC simulation, the quantitative systems pharmacology
metastatic-PDAC simulation with verification, validation, and uncertainty
quantification, the 100,000-patient in silico Phase 3 five-arm PDAC triplicate, and
the end-to-end PDAC digital-twin trial proposals); the H.R. 9510 Verification Before
Generation in Physical AI Oncology Trials Act of 2026 and the national platform and
clinical-trust frameworks; the sixteen-large-language-model code-generation
evaluation against the FDA Adverse Event Reporting System; and the proposed clinical
research protocol v1.0.0 (DOI 10.5281/zenodo.20780121) that this IND operationalizes.

The back matter elements that follow document authorship and contributor
attribution, the open-weights build model used to assemble this application, the
ethical and independence disclosures that bound the entire repository, the data
availability and licensing terms, and the recommended citation. Table 1 records the
provenance of each major content strand so that a reviewer can trace every
quantitative claim in IND Sections 3 through 11 back to its grounding source.

\begin{table}[htbp]
\centering
\caption{Provenance map for the major evidentiary strands of IND v1.0. Each strand
is traced to its grounding source and to the IND sections it supports.}
\begin{tabularx}{\textwidth}{>{\raggedright\arraybackslash}p{3.7cm} Y >{\raggedright\arraybackslash}p{2.5cm}}
\toprule
\textbf{Evidentiary strand} & \textbf{Grounding source and key facts} & \textbf{Supports IND sections} \\
\midrule
Daraxonrasib mechanism and class & Oral RAS(ON) multi-selective (pan-RAS) inhibitor binding GTP-bound mutant RAS in complex with cyclophilin A; once-daily dosing \cite{daraxwhipple} & 3.1, 8.1 \\
Dose-escalation design & 3+3 escalation across DL1 160 mg, DL2 220 mg, DL3 300 mg once daily; 28-day DLT window; RP2D at or below MTD \cite{phase1guidance} & 4, 6.1 \\
Robotic Whipple device & On-premises LLM-directed eight-arm system, 56 degrees of freedom, 640 sensor channels, force caps 3 N per arm and 18 N cumulative \cite{onpremwhippl} & 5, 7 \\
Verification before generation & Ten-gate ACCEPT/BLOCK/ESCALATE suite; LLM as second-opinion advisory oracle only; full autonomy prohibited under 21 CFR \S 312.21(e) \cite{congress9510,clintrustpai} & 5, 10, 11 \\
Perioperative pause-and-restart & ISGPS fistula grade plus serum trough versus 0.5 ng/mL; 32-iteration sweep yielding T+7d 29/32, T+14d 3/32, T+21d 0/32 \cite{paisitedocpk} & 6.1, 9 \\
Computational evidence base & 100,000-patient in silico Phase 3, QSP metastatic-PDAC simulation, and digital-twin trials \cite{chatgpt100kp,qspmetpancre,fdadigtwinpc,pdacdigtwinp} & 8, 9 \\
Regulatory framework & 21 CFR Part 312 and Part 50 Physical AI adaptations; ICH harmonization \cite{cfr312paiuni,cfr050paiuni,ichpaistduni} & 1, 6, 10, 11 \\
\bottomrule
\end{tabularx}
\figcaption{Table 1. Provenance map linking each evidentiary strand to its grounding
reference and to the IND sections it supports.}
\end{table}

\bmhead{Acknowledgments}
This Investigational New Drug Application was assembled with Claude Code (Opus 4.8
Ultracode build, 1M context) as the build model, operating under the single-prompt
mermaid to draft to full to final pipeline described in IND Section 11 and committing
each stage in real time. The author thanks ChatGPT 5.5 (Thinking Extended) and
Gemini 3.1 Pro for research-source assistance and for manual editing support. This
v1.0 application was edited and approved by Kevin Kawchak and added to repository
v4.3.0 under
\href{https://github.com/kevinkawchak/physical-ai-oncology-trials/tree/main/trial-ind}{kevinkawchak/physical-ai-oncology-trials/tree/main/trial-ind}.
The repository contributors are @kevinkawchak, @claude, @google-gemini, and
@openai. No funding body, sponsor, or regulatory authority directed, reviewed, or
endorsed the preparation of this document.

\bmhead{Author and ORCID}
Kevin Kawchak, Chief Executive Officer, ChemicalQDevice, San Diego, California.
ORCID iD:
\href{https://orcid.org/0009-0007-5457-8667}{https://orcid.org/0009-0007-5457-8667}.
Correspondence: \href{mailto:kevink@chemicalqdevice.com}{kevink@chemicalqdevice.com}.
In the proposed study the author would serve as Sponsor-Investigator on behalf of
ChemicalQDevice; the resulting financial interest is disclosed under 21 CFR Part 54
and is mitigated by the independent three-tier oversight structure (Data Safety
Monitoring Board, Independent Safety Monitor, and Physical AI Safety Review
Committee) detailed in the proposed clinical research and additional information
sections.

\bmhead{Ethical Disclosures}
This is an independent research and regulatory-authoring demonstration. No real
patients were enrolled, no protected health information was used, and no
Institutional Review Board approval was sought or obtained. The exemplar case
PAT-PDAC-0001 and all quantitative simulation outputs are synthetic and
illustrative. This document is not an active Investigational New Drug Application or
Investigational Device Exemption, is not a submitted or approved protocol, and does
not constitute medical, surgical, or regulatory advice. It is not endorsed by the
United States Food and Drug Administration, the National Institutes of Health, the
Department of Health and Human Services, any Institutional Review Board, the
International Council for Harmonisation, or any sponsor. Any future first-in-human
conduct of the design described here would require a separately submitted IND and
IDE, a 30-day FDA review period without clinical hold, and prospective IRB approval.

\bmhead{Data Availability}
All application sources, spanning the \texttt{mermaid}, \texttt{draft-ind},
\texttt{full-ind}, and \texttt{final-ind} build stages, are available in the public
GitHub repository and archived on Zenodo under the Creative Commons Attribution 4.0
International (CC BY 4.0) license at DOI
\href{https://doi.org/10.5281/zenodo.xxxxxxxx}{10.5281/zenodo.xxxxxxxx}. The
underlying proposed clinical research protocol v1.0.0 is archived separately at DOI
\href{https://doi.org/10.5281/zenodo.20780121}{10.5281/zenodo.20780121}. No
proprietary, embargoed, or patient-identifiable data are contained in or required to
reproduce this application.

\bmhead{Rights and Permissions}
This article is distributed under the terms of the Creative Commons Attribution 4.0
International License (CC BY 4.0), which permits unrestricted use, distribution, and
reproduction in any medium, provided the original author and source are properly
credited, a link to the Creative Commons license is provided, and any modifications
made are indicated. To view a copy of this license, visit
\url{https://creativecommons.org/licenses/by/4.0/}.

\bmhead{Cite This Article}
Kawchak K. Phase 1 PDAC IND: AI Generation. Zenodo. 2026;
\href{https://doi.org/10.5281/zenodo.xxxxxxxx}{10.5281/zenodo.xxxxxxxx}. IND v1.0,
repository v4.3.0.

\clearpage

% Bibliography (ieeetr); \nocite{*} emits the full assembled reference list.
\phantomsection\addcontentsline{toc}{section}{References}
\begingroup\raggedright\nocite{*}\bibliographystyle{ieeetr}\bibliography{references}\endgroup
